\documentclass[11pt,letterpaper]{article}
% Author: Robert Warren

\usepackage{iftex}
\ifPDFTeX
  \usepackage[utf8]{inputenc}
  \usepackage[T1]{fontenc}
  \usepackage{lmodern}
\fi
\usepackage[margin=1in]{geometry}
\usepackage{hyperref}
\usepackage{RDFTools}

\title{RDF Listing Import Example}
\author{Robert Warren}
\date{\today}

\begin{document}

\pagenumbering{roman}
\listOfRDFfragments

\clearpage
\pagenumbering{arabic}
\maketitle

\section{Imported knowledge graph}

The listing below is numbered, captioned, and cross-referenced like other
document objects. It detects the Turtle serialization from the file extension
and renders the source as a line-numbered listing. See RDF Listing~\ref{rdf:library}
for the linked-data example.

\begin{rdfFragment}[
  source={doc/latex/RDFTools/examples/knowledge-graph.ttl},
  title={Library knowledge graph},
  show-source
]
\caption[Library knowledge graph]{A small linked-data example imported from a Turtle file.}
\label{rdf:library}
\end{rdfFragment}

\section{Direct RDF listing}

When the \texttt{source} option is omitted, the environment captures its body
verbatim. The \texttt{label}, \texttt{caption}, and \texttt{format} options keep
metadata outside the RDF body:

\begin{rdfFragment}[
  title={Inline Turtle example},
  format={Turtle},
  caption={A direct RDF listing},
  label={rdf:inline}
]
@prefix ex: <https://example.org/ns#> .
ex:item ex:label "literal with #, %, _, and {braces}" .
\end{rdfFragment}

The direct listing is numbered and can be referenced as RDF Listing~\ref{rdf:inline}.

\section{Compatibility command}

The original command remains available for simple imports:

\importRDFfragment{doc/latex/RDFTools/examples/knowledge-graph.ttl}

\end{document}
